\section{Transport of slots and a finitely presented central cover}
\label{sec:transport}

We use the fixed finitely presented group $\widehat G$ and its surjection
$\pi:\widehat G\to A_m$ from Section~\ref{sec:lifting}.  Theorem~\ref{thm:polygon-law}
has established the central law for every finite family of aligned polygon
tests.  Our remaining task is to pass from aligned tests to arbitrary
parameterized slots, including several slots in the same track.  We shall
construct their perfect representations in $\widehat G$, prove exact
conjugation formulas, and deduce that $\ker\pi$ is central.
The first obstacle is that an aligned word can fix every point of a
slot while still using its track.  We handle this with private tracks,
first for distinct-track configurations and then for repeated tracks
by proving independence of the routing.  Translation transport is
proved after that independence is available.

We retain the notation
\[
 B_m=\{d=(d_1,\ldots,d_m)\in\Gamma^m:\textstyle\sum_i d_i=0\},
 \qquad t(d)\in\widehat G,
\]
so that $t(d)t(e)=t(d+e)$ exactly.  Write
$\rho_{P,I}:\Alt(I)^\ast\to\widehat G$ for the canonical representation on
an aligned polygon $P$ and a track set $I$, where $|I|\geq5$.  Its image
is perfect.  These maps split exactly over disjoint polygonal partitions,
are compatible under inclusions of track sets, and obey common-translation
covariance on proper track sets, by Theorem~\ref{thm:polygon-law} and
Lemma~\ref{lem:translated-copies}.  We use function composition in which
$ab$ applies $b$ first.

An \emph{aligned word on $J$} means a finite product of elements of the
images of $\rho_{P,I}$, with $I\subseteq J$.  Such words change track labels
but preserve base coordinates after projection.  Their track support is
$J$: balancing tracks used to write translated predicates in the original
generators are not included in $J$.  This convention is legitimate because
Lemma~\ref{lem:translated-copies} proves exact commutation for disjoint
track sets, independently of those choices of balancing words.  In
particular, $t(d)$ commutes with every aligned word on $J$ if $d_i=0$ for
all $i\in J$.

For a finite family with its actual central law, the restriction of
$\pi$ to the subgroup $H$ of conditional copies maps onto the
pointwise group on its nonempty Boolean atoms.  That pointwise group
is a product of alternating groups and acts faithfully on $mX$.
Thus every word in the kernel evaluates to the identity in the actual
assignment model, and the actual central law makes this kernel central.
This supplies the central extensions used below.

Here is the elementary consequence of perfection that will repeatedly
turn a central law into an exact identity.  If $H\to Q$ has central kernel,
$K\leq H$ is perfect, and the image of $b\in H$ centralizes the image of
$K$, then $[b,k]$ is central for every $k\in K$.  The map
$k\mapsto[b,k]$ is therefore a homomorphism from $K$ to an abelian group,
and is trivial.  Thus $b$ centralizes $K$ exactly.  Also, two homomorphisms
from a perfect group to $H$ with the same composition into $Q$ are equal.
Every application below takes place in a specified finite test subgroup
with this central-extension property.

\subsection{Offset frames and distinct tracks}

A parameterized five-slot configuration is data
\begin{equation}
 S=(U;(i_1,u_1),\ldots,(i_5,u_5)),
 \qquad U\subseteq X,\quad i_j\in\{1,\ldots,m\},\quad u_j\in\Gamma,
 \label{eq:slot-configuration}
\end{equation}
where $U$ is a Boolean polygon and the five images of the maps
\[
 s_j:U\longrightarrow mX,\qquad s_j(x)=(i_j,x+u_j)
\]
are pairwise disjoint.  Addition is modulo the coordinate periods, with
the induced action on $X$.  Thus equal track indices are allowed precisely
when the corresponding translated pieces are disjoint.  The ordering
identifies slot permutations with
$E=\Alt(\{1,\ldots,5\})$.  All representations associated to five-slot
configurations will have the same source $E^\ast$.  For an empty $U$ the
representation is the trivial one, and empty pieces in refinements will
be omitted.

For a proper track set $J$ and offsets $a_i\in\Gamma$ for $i\in J$, choose
$f\in B_m$ with $f_i=a_i$ on $J$.  Such an $f$ exists by balancing the sum
on one track outside $J$.  Conjugating aligned representations on $J$ by
$t(f)$ gives their representations in the \emph{offset frame} $a$.  The
choice of $f$ has no effect: if $f'$ is another choice, then $f'-f$ vanishes
on $J$ and $t(f'-f)$ commutes with every aligned word there.  Every finite
central law on $J$ consequently gives the same central law in this frame.

Suppose first that the $i_j$ in~\eqref{eq:slot-configuration} are distinct.
Put $I=\{i_1,\ldots,i_5\}$, let
$\iota_{\boldsymbol i}^\ast:E^\ast\to\Alt(I)^\ast$ be the isomorphism
induced by $j\mapsto i_j$, and choose $f\in B_m$ with $f_{i_j}=u_j$.
Define
\begin{equation}
 \rho_S(s)=t(f)\rho_{U,I}(\iota_{\boldsymbol i}^\ast(s))t(f)^{-1},
 \qquad s\in E^\ast.
 \label{eq:distinct-slot-copy}
\end{equation}
This is independent of the balancing coordinates of $f$, and its
projection is the indicated alternating slot permutation.

\begin{lemma}[Compatibility of frames]\label{lem:slot-frames}
For distinct-track configurations, the representations
\eqref{eq:distinct-slot-copy} have the following properties.
\begin{enumerate}[label=\textup{(\roman*)}]
\item Replacing $U$ by $U+v$ and every $u_j$ by $u_j-v$ leaves the
representation unchanged, with the corresponding reparameterization.
\item If $U=\bigsqcup_{\alpha=1}^r U_\alpha$, then
\begin{equation}
 \rho_S(s)=\prod_{\alpha=1}^r\rho_{S|U_\alpha}(s),
 \label{eq:slot-refinement-distinct}
\end{equation}
and the factor images commute pairwise.
\item For a constant even track permutation $c$, conjugation by its fixed
lift sends $\rho_S$ to the representation of the tuple obtained by
replacing $i_j$ by $c(i_j)$.
\item For every $d\in B_m$,
\begin{equation}
 t(d)\rho_S(s)t(d)^{-1}=\rho_{dS}(s),
 \qquad
 dS=(U;(i_j,u_j+d_{i_j})_{j=1}^5).
 \label{eq:distinct-translation}
\end{equation}
\end{enumerate}
More generally, a representation on a subalphabet computed in an offset
frame agrees with its own frame definition.  Two frames giving the same
offsets on that subalphabet, up to the common reparameterization in
\textup{(i)}, give identical representations there.
\end{lemma}

\begin{proof}
For (i), choose $h\in B_m$ equal to $v$ on $I$.  The common-translation
identity for aligned copies is
$t(h)\rho_{U,I}t(h)^{-1}=\rho_{U+v,I}$.  A balancing vector for the new
offsets, added to $h$, agrees with $f$ on $I$.  Independence of the balance
then proves the assertion.  Part (ii) is the aligned refinement identity
conjugated by $t(f)$.

For (iii), the initial relations give
$c\,t(f)c^{-1}=t(c\cdot f)$ and exact reindexing of the aligned
representations.  These are precisely the data defining the new frame.
For (iv), additivity gives
\[
 t(d)\rho_S(s)t(d)^{-1}
 =t(d+f)\rho_{U,I}(\iota_{\boldsymbol i}^\ast(s))t(d+f)^{-1};
\]
$d+f$ has the required offsets on $I$.  Finally, compatibility under
restriction follows from the corresponding compatibility of aligned
canonical representations and from independence of balances.  This also
proves the assertion about two frames.
\end{proof}

We next prove transport by aligned words.  Private track blocks move
the entire representation away from a word fixing its slots without
changing the commutator.  The lemmas state the spare-track hypotheses
for the indicated support of that word.  Their later applications
compare bounded unions of supports; the final kernel word is treated
one generator at a time.

\begin{lemma}[Fixing a distinct-track tuple]\label{lem:fixed-slot-centralizer}
Let $S$ be a five-slot configuration on distinct tracks, and let $b$ be
an aligned word on a track set $J$.  Suppose that $\pi(b)$ fixes every
point of every parameterized slot of $S$.  Provided the fixed track
reserve contains twenty tracks outside $J$ and the occupied tracks,
together with the balancing tracks used below, $b$ centralizes
$\rho_S(E^\ast)$ exactly.
\end{lemma}

\begin{proof}
For each leg $j$, choose four private tracks
$p_{j,1},\ldots,p_{j,4}$, all distinct and outside $J\cup I$, and put
\[
 V_j=U+u_j,
 \qquad L_j=\{i_j,p_{j,1},\ldots,p_{j,4}\},
 \qquad K_j=\rho_{V_j,L_j}(\Alt(L_j)^\ast).
\]
The projected element $\pi(b)$ fixes the source track pointwise on $V_j$
and fixes every private track.  It therefore commutes with the projection
of $K_j$.  Take the finite family consisting of all tests occurring in a
chosen aligned expression for $b$, the five tests $V_j$, and constants,
on the union of the track sets involved.  Theorem~\ref{thm:polygon-law}
gives a central extension of its pointwise group.  Within that subgroup,
$[b,k]$ is central for $k\in K_j$.  Since $K_j$ is perfect, the
central-commutator argument above proves
\begin{equation}
 [b,K_j]=1\qquad (1\leq j\leq5).
 \label{eq:fixer-private-commutation}
\end{equation}

Choose $a_j\in K_j$ projecting to the cycle
$(i_j\ p_{j,1}\ p_{j,2})$ on $V_j$, and set $a=a_5\cdots a_1$.
There is one common offset frame on the union of the $L_j$: assign offset
$u_j$ to every track in $L_j$.  In this frame, $K_j$ is the aligned
conditional copy on $U$ with track set $L_j$.  Thus $a$ and $\rho_S$ lie
in one conjugate of a subgroup with an aligned central law.  In its
pointwise model, $a$ sends the $j$th slot to
$(p_{j,1},U+u_j)$.  Uniqueness for homomorphisms from $E^\ast$ in this
central extension gives the exact identity
\begin{equation}
 a\rho_S(s)a^{-1}=\rho_T(s),
 \qquad T=(U;(p_{j,1},u_j)_{j=1}^5).
 \label{eq:private-evacuation}
\end{equation}
Here compatibility of subframes in Lemma~\ref{lem:slot-frames} identifies
the two maps with their definitions in~\eqref{eq:distinct-slot-copy}.

The tracks occupied by $T$ are outside $J$.  Write its representation
using a balancing vector $f_T$ which is zero on $J$: prescribe $u_j$ on
$p_{j,1}$ and balance on a further track outside both sets.  The aligned
base copy for $T$ commutes with $b$ by disjoint-track commutation, and
$t(f_T)$ commutes with $b$ because $f_T$ vanishes on $J$.  Consequently
$[b,\rho_T(E^\ast)]=1$.  Equation~\eqref{eq:fixer-private-commutation}
also gives $[b,a]=1$.  Conjugating~\eqref{eq:private-evacuation} back proves
$[b,\rho_S(E^\ast)]=1$.
\end{proof}

\begin{lemma}[Aligned transport between distinct-track tuples]
\label{lem:distinct-slot-transport}
Let $S=(U;(i_j,u_j)_{j=1}^5)$ and
$T=(U;(k_j,u_j)_{j=1}^5)$ be configurations, each with distinct track
indices.  If an aligned word $b$ sends the $j$th parameterized slot of
$S$ to the $j$th slot of $T$, then
\begin{equation}
 b\rho_S(s)b^{-1}=\rho_T(s)\qquad(s\in E^\ast),
 \label{eq:distinct-aligned-transport}
\end{equation}
whenever the fixed reserve supplies the tracks of
Lemma~\ref{lem:fixed-slot-centralizer} for $b$ and an even constant
reindexing of the two tuples.
\end{lemma}

\begin{proof}
Extend the bijection $i_j\mapsto k_j$ to a permutation of the union of
the two five-element sets.  If it is odd, multiply by a transposition on
two additional tracks.  The resulting even constant permutation $c$
has support at most twelve and sends $S$ to $T$ parameterwise.
The aligned word $c^{-1}b$ fixes $S$ pointwise.  Apply
Lemma~\ref{lem:fixed-slot-centralizer}, and then the constant-reindexing
identity of Lemma~\ref{lem:slot-frames}.
\end{proof}

\subsection{Repeated tracks and independence of routing}

We now construct the representation of~\eqref{eq:slot-configuration}
without a distinctness assumption on the $i_j$.  An \emph{admissible
routing} of $S$ is an aligned word $r$ on an alphabet of at most
$25$ tracks containing all source and terminal track labels, whose
projected action sends its five slots parameterwise to a tuple $T$ on
distinct tracks.  Since an aligned word preserves base
coordinates, $T$ has the form $(U;(k_j,u_j)_{j=1}^5)$.  A routing need not
act trivially outside the specified slots; this is useful when the same
routing is applied to a subdivision of $U$.

Such a routing always exists.  Choose four private tracks for each leg,
all outside the occupied source tracks and all different.  On
$V_j=U+u_j$ choose an element of the perfect aligned copy on
\[
 L_j=\{i_j,p_{j,1},p_{j,2},p_{j,3},p_{j,4}\}
\]
projecting to $(i_j\ p_{j,1}\ p_{j,2})$.  Multiply these five elements
to obtain $r$.  Each factor sends its own slot to $p_{j,1}$.  It fixes
every other source slot: different source tracks are not in its private
block, and slots sharing $i_j$ have disjoint Boolean regions by
hypothesis.  Private destinations of different legs are distinct.
The total support is at most $5+5\cdot4=25$.

For an admissible routing $r:S\to T$, define provisionally
\begin{equation}
 \rho_S^{\,r}(s)=r^{-1}\rho_T(s)r.
 \label{eq:routed-copy}
\end{equation}
The next lemma proves equality of these maps in $\widehat G$ itself.

\begin{lemma}[Independence of routing]\label{lem:routing-independence}
For every five-slot configuration $S$, the homomorphism
\eqref{eq:routed-copy} is independent of its admissible routing.  Denote
it by $\rho_S:E^\ast\to\widehat G$.  It agrees with
\eqref{eq:distinct-slot-copy} when the tracks are distinct.  For every
finite polygonal partition $U=\bigsqcup_\alpha U_\alpha$,
\begin{equation}
 \rho_S(s)=\prod_\alpha\rho_{S|U_\alpha}(s),
 \label{eq:slot-refinement}
\end{equation}
with pairwise commuting factor images.  Common reparameterization of
$U$ and the $u_j$ leaves $\rho_S$ unchanged.
\end{lemma}

\begin{proof}
Let $r_1:S\to T_1$ and $r_2:S\to T_2$ be two admissible routings.  The
aligned word $b=r_2r_1^{-1}$ sends $T_1$ to $T_2$ parameterwise.  Both
terminal tuples have distinct tracks, so Lemma~\ref{lem:distinct-slot-transport}
applies and gives
\[
 r_2r_1^{-1}\rho_{T_1}(s)r_1r_2^{-1}=\rho_{T_2}(s).
\]
After multiplying by $r_2^{-1}$ and $r_2$, this is precisely
$\rho_S^{\,r_1}(s)=\rho_S^{\,r_2}(s)$.  A distinct-track tuple admits the
identity routing, proving agreement with its earlier definition.

Fix one routing $r:S\to T$.  It is also an admissible routing of every
$S|U_\alpha$ to $T|U_\alpha$: its support bound has not changed.  Apply
Lemma~\ref{lem:slot-frames}(ii) to $T$, and conjugate the identity and its
commutation assertions by $r^{-1}$.  This proves~\eqref{eq:slot-refinement}.
For common reparameterization, use the same routing, which depends only
on the actual parameterized slot maps after their domain identification,
and apply Lemma~\ref{lem:slot-frames}(i) at its distinct-track target.
\end{proof}

Routing independence has now been established using aligned words only.
In particular, it is available before any claim that a translation
transports a repeated-track configuration.

\begin{lemma}[Aligned transport of arbitrary tuples]
\label{lem:aligned-slot-transport}
Let $b$ be an aligned word supported on at most five tracks.  Suppose
that it sends a five-slot configuration $S$ parameterwise to a
configuration $T$ with constant track labels on each leg.  Then
\begin{equation}
 b\rho_S(s)b^{-1}=\rho_T(s)\qquad(s\in E^\ast).
 \label{eq:aligned-slot-transport}
\end{equation}
If track labels vary within $U$, the same assertion holds piecewise on
a finite polygonal partition of $U$, with the product interpretation
of~\eqref{eq:slot-refinement}.
\end{lemma}

\begin{proof}
Choose admissible routings $r_S:S\to S'$ and $r_T:T\to T'$.  The word
$r_Tbr_S^{-1}$ is aligned and sends the distinct-track tuple $S'$ to
$T'$.  Lemma~\ref{lem:distinct-slot-transport} gives
\[
 (r_Tbr_S^{-1})\rho_{S'}(s)(r_Tbr_S^{-1})^{-1}=\rho_{T'}(s).
\]
Conjugate by $r_T^{-1}$ and use the definitions of $\rho_S$ and $\rho_T$.
The support of the transition word is at most $25+5+25=55$, so this
application uses a fixed reserve.

For the last assertion, the finitely many polygon tests in $b$ induce
a finite partition on each map $x\mapsto x+u_j$.  Their common refinement
makes all five terminal track labels constant.  Apply the proved
identity on every part and multiply, using~\eqref{eq:slot-refinement}.
\end{proof}

\subsection{Balanced translations}

The representations on repeated tracks are now independent of routing.
We next use that independence to prove translation covariance.  The
auxiliary translation below shifts each chosen routing block uniformly;
a second routing handles the difference from the desired translation.

\begin{lemma}[Translation transport]\label{lem:translation-slot-transport}
For every $d\in B_m$ and every five-slot configuration $S$,
\begin{equation}
 t(d)\rho_S(s)t(d)^{-1}=\rho_{dS}(s),
 \qquad s\in E^\ast,
 \label{eq:translation-slot-transport}
\end{equation}
where $dS=(U;(i_j,u_j+d_{i_j})_{j=1}^5)$.
\end{lemma}

\begin{proof}
It suffices to treat the fixed generators of $B_m$, each supported on
two tracks, and their inverses.  The general statement then follows by
composition and additivity of $t$.

Choose the private-block routing $r:S\to T$ used in the existence proof
above.  Its $j$th factor lies in $\rho_{V_j,L_j}(\Alt(L_j)^\ast)$, where
$V_j=U+u_j$.  Choose $d'\in B_m$ as follows:
\begin{equation}
 d'_i=d_i\quad\text{on every occupied source track }i,
 \qquad d'_{p_{j,a}}=d_{i_j}\quad(1\leq j\leq5,\ 1\leq a\leq4).
 \label{eq:blockwise-translation}
\end{equation}
Put $d'=0$ on the other tracks except for one unused track, on which
we balance the sum.  Repeated source labels give consistent prescriptions
because they prescribe the same $d_i$.  Thus $d'$ is constant on each
$L_j$, and has support at most $26$.  Set
\[
 \delta=d-d'.
\]
The translation $\delta$ vanishes on all occupied source tracks and has
bounded support (at most $28$ is sufficient).

We first show
\begin{equation}
 [t(\delta),\rho_S(E^\ast)]=1.
 \label{eq:delta-centralizes}
\end{equation}
Choose another private-block routing $q:S\to T_0$, with all its private
tracks outside the support of $\delta$.  Every track in each of its
five blocks then has $\delta_i=0$.  Outside-support translation
commutation, from Lemma~\ref{lem:translated-copies}, gives
$[t(\delta),q]=1$ factor by factor.  The tuple $T_0$ has distinct tracks,
all outside the support of $\delta$.  Write
\[
 \rho_{T_0}(s)=t(f)\rho_{U,I_0}
       (\iota_{\boldsymbol k}^\ast(s))t(f)^{-1}.
\]
The element $t(\delta)$ commutes with $t(f)$ by additivity and with the
aligned representation on $I_0$ because it is zero on $I_0$.
It therefore centralizes $\rho_{T_0}(E^\ast)$.  By the already proved
routing independence,
$\rho_S=q^{-1}\rho_{T_0}q$, proving~\eqref{eq:delta-centralizes}.
This argument uses only the distinct-track frame formula and
outside-support commutation.

Next conjugate the original routing by $t(d')$.  Since $d'$ is constant
on $L_j$, the common-translation law for aligned polygon representations
replaces its $j$th conditional factor on $V_j$ by the same factor on
$V_j+d_{i_j}$.  Hence
\[
 r'=t(d')rt(d')^{-1}
\]
is an aligned word on the same at-most-$25$ tracks.  It routes $d'S$
to $d'T$, preserving base coordinates at this new configuration.  It is
therefore an admissible routing.  The distinct-track formula
\eqref{eq:distinct-translation} gives
\[
 \begin{split}
 t(d')\rho_S(s)t(d')^{-1}
 &=r'^{-1}\bigl(t(d')\rho_T(s)t(d')^{-1}\bigr)r'\\
 &=r'^{-1}\rho_{d'T}(s)r'
 =\rho_{d'S}(s).
 \end{split}
\]
Finally, $d'S=dS$ because $d'=d$ on the occupied source tracks, while
$t(d)=t(d')t(\delta)$.  Equation~\eqref{eq:delta-centralizes} now proves
\eqref{eq:translation-slot-transport}.
\end{proof}

\paragraph{The fixed number of tracks.}
We record why the preceding constructions fit the single choice of $m$
made before the presentation was fixed.  One private routing uses at most
$25$ tracks.  Comparing two routings uses at most $50$, and inserting one
five-track aligned generator uses at most $55$.  The even constant
reindexing in Lemma~\ref{lem:distinct-slot-transport} uses at most twelve
tracks.  The fixing-slot argument adds twenty private tracks and finitely
many balancing tracks.  The translation argument uses the source tracks,
two banks of twenty private tracks, the at-most-two-track support of its
generator, and a balancing track; its temporary difference translation
has support at most $28$.  These are absolute finite bounds, independent
of polygons, offsets, the number of refinement pieces, and word length.
Let $M_{\mathrm{tr}}$ be the maximum of the total track requirements
of these finitely many operations, including the auxiliary tracks in
each.  Taking $m>M_{\mathrm{tr}}$ in addition to the earlier support
requirements makes every comparison available.

A refined piece uses the same routing as its parent.  Different pieces
are processed successively, and successive word letters may reuse the
private banks because Lemma~\ref{lem:routing-independence} identifies all
choices exactly.  No bank is assigned permanently to a predicate, a
refinement level, or a letter of a later word.  The only relations used
here are the initial additive and reindexing relations, outside-support
commutation, and the polygon laws already deduced in the fixed presented
group.  In particular, this reserve does not require any new relations
depending on the configuration being transported.

\subsection{Generation and centrality of the kernel}

Let $N\leq\widehat G$ be the subgroup generated by
$\rho_S(E^\ast)$ over all parameterized five-slot configurations $S$.
We first check that these copies generate the presented group itself.
This is stronger than their generation after projection and is necessary
for the kernel argument.

\begin{lemma}[Generation by slot copies]\label{lem:slot-generation}
The subgroup $N$ is all of $\widehat G$.
\end{lemma}

\begin{proof}
Each of the initial conditional alternating track groups is generated
by its subgroups on five tracks.  The canonical representations on those
five tracks agree exactly with their initially specified lifts, by the
finite initial tables and uniqueness for perfect sources.  Thus $N$
contains all conditional track-permutation generators, including the
constant ones.

It remains to include the balanced translation generators.  Let $v$ be
one of the two fixed generators of $\Gamma$, and choose four distinct
tracks $i,j,k,h$.  Put
\[
 c=(i\ j\ k),\qquad a=v e_i-v e_h\in B_m,
\]
where $e_i$ denotes placement in the $i$th track coordinate.  We use the
same symbol $c$ for the specified constant lift.  The additive and
constant-reindexing relations give
\begin{equation}
 [c,t(a)]=t(c\cdot a-a)=t(v e_j-v e_i).
 \label{eq:translation-as-slot-product}
\end{equation}
The commutator convention is $[x,y]=xyx^{-1}y^{-1}$.  On the other hand,
\[
 [c,t(a)]=c\,\bigl(t(a)c^{-1}t(a)^{-1}\bigr).
\]
The first factor belongs to a constant five-track representation, after
adjoining two track labels to the three moved by $c$.  The second belongs
to its offset-frame conjugate and hence, by
\eqref{eq:distinct-slot-copy}, to another five-slot representation.
Both factors are in $N$, so $t(v e_j-v e_i)\in N$.

Such differences generate $B_m$.  Only the fixed basis differences
are needed for the finite generator list, and the finitely many instances
of~\eqref{eq:translation-as-slot-product} are among the initial relations
of Section~\ref{sec:lifting} (they also follow from its additive and
reindexing relations).  This proves $N=\widehat G$ without assuming that
any conjugating translation already belongs to $N$.
\end{proof}

\begin{theorem}[A finitely presented central cover]\label{thm:central-cover}
For all sufficiently large fixed finite $m$, there is a finitely
presented group $\widehat G$ and a surjective homomorphism
$\pi:\widehat G\to A_m$ whose kernel is central.
\end{theorem}

\begin{proof}
Choose $m$ beyond the fixed track requirements of the lifting and
propagation arguments and $M_{\mathrm{tr}}$.  Choose the
finite initial relation menu in the order specified in
Sections~\ref{sec:lifting} and~\ref{sec:propagation}.  This produces the
finitely presented group $\widehat G$ and the surjection $\pi$; the
transport lemmas have all been proved inside this group.

Take $w\in\ker\pi$, and write it as a finite word in the chosen
generators and their inverses.  Each conditional permutation letter
can be written, using its initial finite group table, as a product of
letters supported on five tracks.  Thus we may regard every letter
as either a five-track aligned element or a fixed balanced translation
generator or its inverse.  Write the resulting word as $w=g_n\cdots g_1$
and put $w_k=g_k\cdots g_1$, so that the letters act in the order
$g_1,\ldots,g_n$.  Fix a configuration
$S=(U;(i_j,u_j)_{j=1}^5)$.

There is a finite polygonal partition $U=\bigsqcup_\alpha U_\alpha$
such that, along every $w_k$ and for each of the five
legs, its image on $U_\alpha$ has one fixed track label and one fixed
translation offset.  To construct the partition, pull back the finitely
many continuity pieces of each next letter along the five current
partial translations and intersect with the pieces already obtained.
The Boolean algebra is closed under these operations.  Induction on
the finite word length therefore gives a finite partition; there is no
claimed bound on its number of pieces.

On a nonempty $U_\alpha$, follow the resulting five-slot tuple through
successive letters.  The images remain disjoint because every projected
letter is a bijection.  Lemma~\ref{lem:aligned-slot-transport} applies
to aligned letters and Lemma~\ref{lem:translation-slot-transport} to
translation letters.  Composing their exact conjugation identities
gives
\begin{equation}
 w\rho_{S|U_\alpha}(s)w^{-1}
   =\rho_{\pi(w)(S|U_\alpha)}(s).
 \label{eq:word-slot-transport}
\end{equation}
Since $\pi(w)=1$, the terminal tuple equals the initial tuple
parameterwise.  In particular the terminal track labels are the same;
the terminal offsets are the same in $\Gamma$, since the translation
action on $X$ is free.  The right side of~\eqref{eq:word-slot-transport}
is consequently $\rho_{S|U_\alpha}(s)$.

Multiply these identities over $\alpha$ and use the exact refinement
formula~\eqref{eq:slot-refinement}.  We obtain
$[w,\rho_S(E^\ast)]=1$ for every configuration $S$.
Lemma~\ref{lem:slot-generation} says that these images generate
$\widehat G$, so $w\in Z(\widehat G)$.  Since $w$ was arbitrary,
$\ker\pi\subseteq Z(\widehat G)$, as required.
\end{proof}
